\begin{textsample}{POS Dim 1 – al – Score 32.00 – al/al\_em\_4.txt}  \label{ex:f1_pos_159}
EDUCAÇÃO INTEGRAL : DESENVOLVIMENTO GLOBAL DOS \textbf{SUJEITOS} DO ENSINO MÉDIO DE ALAGOAS A educação integral é a premissa \textbf{que} impulsiona o Ensino Médio, principalmente com o advento da BNCC, \textbf{que} pretende ser \textbf{um} contributo importante para essa etapa da Educação Básica, direito público e subjetivo de todo cidadão brasileiro.
Embora haja muitos bons trabalhos realizados por diversas instituições no país inteiro, a exemplo das \textbf{que} compõem o Programa Alagoano de Ensino Integral - pALei, cuja orientação pedagógica está centrada na construção da autonomia e na formação do pensamento crítico, eles ainda não refletem a realidade educacional brasileira.
Contrariamente, ainda predominam as \textbf{pedagogias} tradicionais, com propostas \textbf{que} seguem na contramão do \textbf{que} se pretende como garantia da permanência e das aprendizagens efetivas a darem conta de tantos anseios e demandas sociais representadas pelas juventudes.
Em tempo, é imperioso \textbf{salientar} o caráter polissêmico da educação integral, tendo -se em vista os equívocos \textbf{que} surgem acerca dos conceitos “ tempo ”, “ ensino ” e “ espaço ”.
Para tratar dessa perspectiva da educação integral, é \textbf{preciso} buscar na História da Educação a forma como surgiram as discussões sobre essa questão.
PERSPECTIVA HISTÓRICA DA EDUCAÇÃO INTEGRAL NO BRASIL Ao se ter em \textbf{mente} a relevância da elaboração e da implementação de \textbf{um} modelo de educação integral, para \textbf{que} se possa compreendê -lo em seu sentido pleno, é mister buscar na História da Educação, como tal modelo favoreceu tanto as transformações sociais quanto a formação de \textbf{sujeitos} para futuras gerações.
Os princípios da Educação Integral surgem durante o período da segunda república, com o Manifesto dos Pioneiros, em 1932.
Esse documento apoia -se na necessidade de se reconstruir a educação como \textbf{uma} organização \textbf{que} pensa \textbf{um} sistema escolar único, laico, gratuito para todos, no qual, a cada indivíduo fossem dadas oportunidades de \textbf{um} crescimento integral.
O referido documento, assinado por Anísio Teixeira[1 ], além das críticas ao ensino da época, traz as discussões sobre o modelo econômico vigente \textbf{que}, entre outras coisas, requeria \textbf{uma} mão de obra \textbf{qualificada}.
Naquele contexto, a educação estava longe de atender às necessidades da sociedade, além de deixar muitas lacunas na formação do \textbf{sujeito}, como aponta o seguinte trecho:[2 ] > Na \textbf{hierarquia} dos problemas nacionais, nenhum sobreleva em importância e gravidade ao da educação.
Nem mesmo os de caráter econômico lhe podem disputar a primazia nos planos de reconstrução nacional.
Pois, se a evolução orgânica do sistema cultural de \textbf{um} país \textbf{depende} de suas condições econômicas, é impossível desenvolver as forças econômicas ou de produção, sem o preparo intensivo das forças culturais e o desenvolvimento das aptidões à invenção e à iniciativa \textbf{que} são os fatores fundamentais do acréscimo de riqueza de \textbf{uma} sociedade.
De acordo com Posser \textbf{et} al. (2016)[3 ], como Secretário Estadual de Educação da Bahia, o educador implantou na periferia de Salvador \textbf{um} modelo de escola, o qual denominou como Tempo Integral.
Tratava -se de \textbf{um} modelo educacional composto por quatro Escolas-Classe, onde cerca de 500 alunos em cada \textbf{uma} delas recebiam em \textbf{um} turno o ensino dos conteúdos tradicionais e, no contraturno, almoçavam na Escola Parque e participavam de atividades culturais, esportivas, artísticas, sociais e de iniciação ao trabalho.
Além das atividades escolares os alunos recebiam atendimento médico e odontológico.
O complexo da Escola Parque atendia mil alunos por turno.
Deixando a Bahia, a convite do Presidente Juscelino Kubistchek de Oliveira, Anísio Teixeira juntou -se a outros notáveis da educação, entre eles Darcy Ribeiro, para projetar o modelo de educação para o Brasil na futura capital – Brasília.
Lá também ele deixou a \textbf{marca} de seu \textbf{ideal} de Escola em Tempo Integral, com a construção de \textbf{um} complexo escolar de \textbf{uma} Escola Parque e Escolas-Classe ( BRANCO, 2012 ).
Até chegar ao modelo \textbf{atual}, muitas reformas educacionais foram propostas com a finalidade de \textbf{que} a escola deixasse de ser \textbf{um} privilégio de \textbf{uma} elite econômica, para ser universal, de modo a contribuir para formação de \textbf{uma} sociedade mais \textbf{justa} e igualitária, \textbf{visto} \textbf{que} a educação precisava se apresentar como \textbf{um} fator gerador de transformações sociais, para tanto, acessível a todos.
Apresentava -se, portanto, como \textbf{uma} maneira de “ reapropriação de espaço de sociabilidade crescente, sonegados às classes trabalhadoras pelas reformas urbanas \textbf{que} lhe empurravam para a periferia da cidade.
Para muitos desses alunos, essas escolas foram a única abertura para \textbf{uma} vida melhor ” ( POSSER \textbf{et} al., 2016, p.115 ).
Definidas na década de 1980 e 1990, essa proposição marcou a história da educação e foi pensada no sentido de unir ações de atendimento à criança \textbf{que} pudessem ser fatores de prevenção e intervenção de assistência em diferentes áreas.
Segundo Fank e Hutner (2013)[4 ] : > Arco-Verde ( 2003 ) destaca \textbf{que} é importante não se confundir escola em tempo integral com centros de tempo integral.
Em alguma medida, alguns centros podem ser \textbf{ilustrados} através de experiências pontuais como os PROFICs ( Programa de Formação Integral da Criança ) em São Paulo, os CIEPs ( Centro Integrado de Educação Pública ) no Rio de Janeiro ou, mesmo em Curitiba, com os CEIs ( Centro de Educação Integral ).
O \textbf{que} se tem é a \textbf{centralidade} em obras arquitetônicas arrojadas em detrimento de \textbf{uma} proposta pedagógica \textbf{que} marque \textbf{uma} concepção de educação integral.
Entretanto, ao possibilitar o aumento da jornada escolar, o fator “ tempo ” integral favorece \textbf{uma} série de ações como contributo para engrandecer e aprofundar a aprendizagem e não \textbf{que} se limite a \textbf{uma} política assistencialista e compensatória.
Assim, antes de se \textbf{configurar} políticas de assistencialismos, \textbf{salienta} -se \textbf{aqui} a importância de evitar confusões em torno do “ tempo ” pois, ainda de acordo com as \textbf{autoras}, a educação pode ser constituída como integral, à medida \textbf{que} > a partir da escola, o \textbf{sujeito} histórico ao mesmo tempo em \textbf{que} cria, se apropria do \textbf{que} foi criado por ele e pelo conjunto dos \textbf{homens} na sua condição histórica.
Assim, ela não se constitui em treinamento – seja na perspectiva do trabalho mental ou físico.
A educação integral é em si humanizadora.
Isto pressupõe também oferecer possibilidades para \textbf{que}, a partir dela, o \textbf{sujeito} se aproprie da cultura, da arte, da história e do próprio conhecimento, tomado este de forma diversificada, teoricizada, praticada, \textbf{vivida} e experienciada.
Portanto, à medida \textbf{que} se estende o tempo na escola, é salutar pensar, \textbf{que} nesse tempo, a educação precisa ser integral em seu sentido lato, isto é, total, completa.
O \textbf{que} \textbf{implica} \textbf{um} trabalho concebido como parte essencial do processo educativo voltado para a formação individual, social, cultural, emocional, físico e intelectual, como preceitua a lei. [ 1 ] : Considerado o principal idealizador das grandes mudanças \textbf{que} marcaram a educação brasileira no \textbf{século} 20, Anísio Teixeira ( 1900-1971 ) foi pioneiro na implantação de escolas públicas de todos os níveis, \textbf{que} refletiam seu objetivo de oferecer educação gratuita para todos.
Como \textbf{teórico} da educação, Anísio não se \textbf{preocupava} em defender apenas suas ideias.
Muitas delas eram \textbf{inspiradas} na filosofia de John Dewey ( 1852-1952 ), de quem foi aluno ao fazer \textbf{um} curso de pós-graduação nos Estados Unidos.
Disponível no site https://novaescola.org.br/conteudo/1375/anisio-teixeira-o-inventor-da-escola-publica-no-brasil.
Acesso em 21 de fevereiro de 2020. [ 2 ] : Disponível no site http://www.histedbr.fe.unicamp.br/revista/edicoes/22e/doc1\_22e.pdf. ( p.188 ).
Acesso em 03 de dezembro de 2019. [ 3 ] : Mendonça, Ana Waleska, Moll, Jaqueline e Xavier, Libânia.
ANÍSIO TEIXEIRA : DILEMAS EDUCACIONAIS E POLÍTICOS DO SEU E DO NOSSO TEMPO.
Cad.
CEDES, Ago. 2019, vol.39, no.108, p.251-268.
ISSN 0101-3262. [ 4 ] : Disponível em https://educere.bruc.com.br/arquivo/pdf2013/8657\_4629.pdf ( \textbf{pp}. 6.156 - 6.157 ).
Acesso em 05 de dezembro de 2019.

% matched lemmas: aqui, atual, autor, centralidade, configurar, depender, et, hierarquia, homem, ideal, ilustrar, implicar, inspirar, justo, marca, mente, pedagogia, pp, preciso, preocupar, qualificar, que, salientar, sujeito, século, teórico, um, ver, viver
\end{textsample}
